\section{Introduction}
\label{sec:intro}

Benchmarks for machine-learning agents have settled into a shape. You give the
agent a dataset and a metric, and you see how good a model it can fit. MLE-bench
draws its tasks from Kaggle competitions; MLAgentBench, DSBench,
ScienceAgentBench and their relatives vary the domain and the scaffolding but
keep the shape.\footnote{Citations are collected in
Section~\ref{sec:related}.} In every case the data exists before the agent
starts, and it is the same data for every agent. That is exactly what makes
those benchmarks clean comparisons.

It is also what makes them unlike most of computational science. There, the
expensive step is usually \emph{getting} the data --- running the simulation,
performing the experiment, doing the solve --- and the decision that separates a
good researcher from a mediocre one is \emph{where to spend the next unit of
budget}. An agent that is excellent at fitting a dataset it was handed may be
useless at deciding which experiments are worth running. Nothing in the current
crop of benchmarks would tell us either way.

This paper introduces a benchmark that puts that decision back. The agent is
given a physics solver, a five-dimensional space of inputs, and a hard budget of
solver calls. It must work out how to drive the solver, decide which points in
the space to spend its budget on, train a model on what it collected, and hand
over a predictor --- which we then score, on a set of cases it has never seen,
with machinery it does not control (Figure~\ref{fig:loop}).

\begin{figure}[t]
\centering
\resizebox{\textwidth}{!}{%
\begin{tikzpicture}[
  font=\small,
  box/.style={draw=black!55, rounded corners=2pt, align=center, inner sep=4pt,
              minimum height=9mm, text width=30mm},
  ag/.style={box, fill=blue!6,   draw=blue!60},
  sv/.style={box, fill=orange!8, draw=orange!70},
  ev/.style={box, fill=green!8,  draw=green!45!black},
  dl/.style={box, fill=black!4},
  fl/.style={-{Stealth[length=2mm]}, draw=black!70, thick},
  lb/.style={font=\scriptsize\itshape, text=black!60},
  hd/.style={font=\scriptsize\bfseries, text=black!70},
]

% ---- left: inside the evaluated cell ----
\node[dl] (task)  at (0,0)      {frozen task spec\\[1pt]{\scriptsize 5 parameters $\cdot$ solve budget\\ process gates}};
\node[ag] (agent) at (0,-2.1)   {agent harness\\[1pt]{\scriptsize \Ltwo{} library-assisted\\ \Lthree{} from scratch}};
\node[sv] (solv)  at (5.4,-2.1) {TokaMaker\\[1pt]{\scriptsize \GS{} FE solver}};
\node[dl] (deliv) at (0,-4.6)   {deliverable\\[1pt]{\scriptsize \code{report.json}\\ \code{predict.py}}};

\draw[fl] (task) -- (agent);
\draw[fl] ([yshift=3.5mm]agent.east) -- ([yshift=3.5mm]solv.west)
      node[midway,above,lb]{params};
\draw[fl] ([yshift=-3.5mm]solv.west) -- ([yshift=-3.5mm]agent.east)
      node[midway,below,lb]{$\psi$ or failure};
\draw[fl] (agent) -- (deliv);

\node[lb, text width=36mm, align=center] at (5.4,-3.2)
     {$\le 450$ solves, 3 adaptive rounds};

% ---- separator ----
\draw[dashed, black!45] (7.7,1.0) -- (7.7,-5.5);
\node[hd, anchor=south] at (2.7,0.85) {inside the evaluated cell};
\node[hd, anchor=south] at (10.9,0.85) {independent evaluation};

% ---- right: evaluation the agent never sees ----
\node[ev] (gates) at (10.9,-0.4) {contract gates\\[1pt]{\scriptsize did it deliver?}};
\node[ev] (score) at (10.9,-2.5) {frozen test set\\[1pt]{\scriptsize 60 held-out solves\\ relative $L_2$}};
\node[ev] (diag)  at (10.9,-4.6) {diagnostics\\[1pt]{\scriptsize valid? honest?}};

\coordinate (bus) at (8.6,-4.6);
\draw[fl] (deliv.east) -- (bus);
\draw[fl] (bus) -- (diag.west);
\draw[fl] (bus) |- (score.west);
\draw[fl] (bus) |- (gates.west);

\end{tikzpicture}%
}
\caption{The benchmark loop. Everything inside the dashed line is the agent's
own business: it is handed a frozen task specification and a solver, and it must
decide which equilibria to compute, spending a fixed budget across an initial
design and up to three adaptive rounds. What it hands over is a small set of
files. Everything to the right of the dashed line happens outside the agent's
reach --- the files are checked mechanically, the predictor is run on sixty
held-out cases the agent never sees, and the result is put through diagnostics
that ask whether the field is physically meaningful and whether the agent's own
report of its performance was true. Every condition in this paper, including
the two that use no language model at all, passes through that identical
right-hand column.}
\label{fig:loop}
\end{figure}

\paragraph{Why this particular solver.}
The ground truth comes from TokaMaker~\citep{hansen2024tokamaker}, a
finite-element solver for the equilibrium of a magnetically confined fusion
plasma. Section~\ref{sec:physics} explains the physics from scratch; the
properties that make it a good instrument for measuring agents are these. It is
\emph{genuinely expensive} --- seconds per solve, so a budget bites. It
\emph{fails}, and does so as a function of the inputs, so the feasible region
has to be discovered rather than assumed. Its output is
\emph{high-dimensional} (a field on a $96\times64$ grid) while its input is
five numbers, so dimensionality reduction is a real design decision rather than
a formality. It is \emph{unforgiving about physics}: the quantity being
predicted is a signed physical field in webers, undefined outside the plasma, so
an agent that quietly normalises it, or fills in the region where it does not
exist, produces something that looks plausible and is wrong. And the answers are
\emph{not memorisable} --- these particular equilibria are not on the internet.

\paragraph{What we found.}
Three things, in increasing order of how much they surprised us.

First, with the language model held fixed across every condition, \emph{which
agent framework you use matters more than how much help you give it}. The
per-cell medians span a factor of $8.8$ across four substrates, against a factor
of $1.7$ between the two levels of scaffolding.

Second, \emph{the scaffolding helped in the wrong direction}. Agents allowed to
import a proven, tested pipeline library did worse than agents made to write
everything themselves --- in all four substrates, with no exceptions. A
deterministic analysis of the code they wrote (Section~\ref{sec:agentlogic})
says why, and it is not that the library is bad: given the library, all four
substrates adopt its default model and the acquisition strategies its
documentation names; given nothing, all four independently find a different
method family, and two of them beat the library's own score on the same test
set.

Third, and this is the finding we think matters beyond this benchmark:
\textbf{31\% of the runs that passed every machine-checkable gate were
physically invalid}. Not marginal --- invalid. One of them scored better than
most of the corpus on the error metric while emitting a magnetic field
throughout the vacuum region outside the plasma, where the field is not defined
at all (Figure~\ref{fig:pva}, second row). The contract that certified it is the
standard instrument in this literature: check the files exist, check they parse,
check the predictor runs and returns the right shape. It cannot see this class
of failure, and neither can the headline error number.

\paragraph{Contributions.}
\begin{enumerate}
\item \textbf{A closed-loop scientific-ML benchmark} in which the agent buys its
  own training data from an expensive solver under a fixed budget, with a
  machine-checkable deliverable, independent scoring on a frozen set, and a
  four-rung ladder of scaffolding from a zero-LLM scripted policy to
  from-scratch --- every rung inside the identical contract
  (Section~\ref{sec:benchmark}).
\item \textbf{A model-controlled comparison of four agent frameworks.} Every
  cell runs the same model, so ``which framework'' is not confounded with
  ``which model'' --- a confound we had in our own earlier campaigns
  (Section~\ref{sec:results}).
\item \textbf{A verification layer} that separates delivery, physical validity
  and self-report honesty, and finds all three come apart
  (Section~\ref{sec:verification}).
\item \textbf{A deterministic, LLM-free account of what each agent actually
  built} --- its chain of methods, its per-round decision logic, and whether its
  code agrees with its own prose. This is what makes the scaffolding result
  interpretable rather than merely surprising (Section~\ref{sec:agentlogic}).
\end{enumerate}
